% ============ 02. Resumen ejecutivo ============

\section{Resumen ejecutivo}

\subsection{Objeto de la auditoría}

Este documento presenta la auditoría integral del repositorio
\texttt{Observatorio\_Ministerio\_de\_Ciencias\_Grupo8}
(\url{https://github.com/ustadistica/Observatorio_Ministerio_de_Ciencias_Grupo8}),
que materializa un \textbf{observatorio crítico del sistema de reconocimiento de
investigadores del Ministerio de Ciencia, Tecnología e Innovación de Colombia
(MinCiencias)}. El análisis cubre 236 archivos (51 de código), 133 commits
(2025-08-24 a 2026-05-14), 12 autores y un pipeline que procesa 77\,237 registros
del padrón de investigadores (6 convocatorias, 2013--2021) y 3\,166\,629 filas de
producción de grupos de investigación.

\subsection{Veredicto general}

\begin{tcolorbox}[hallazgo]
El repositorio cumple su objetivo declarado: consolidar, validar y analizar el
padrón de investigadores reconocidos y producir un observatorio con 14 hallazgos
críticos documentados, dashboard de 7 capítulos, informe LaTeX y presentación.
La \textbf{arquitectura y la lógica estadística de nivel descriptivo-exploratorio
son sólidas}; la \textbf{deuda técnica} (CI/CD rota, Dockerfile desalineado,
codificación mojibake en el CSV canónico, tests mínimos, legacy R no reproducible,
versionado de datos incompleto) compromete la \textbf{reproducibilidad end-to-end},
pero no la validez de los hallazgos ya generados.
\end{tcolorbox}

\subsection{Principales fortalezas verificadas}

\begin{itemize}
  \item \celdaok{Documentación de nivel profesional}: README exhaustivo, catálogo
        YAML con diccionario de datos y llaves de cruce, decisiones metodológicas
        documentadas en el propio código (p.~ej. tratamiento de atípicos en
        \texttt{src/analisis/calidad.py}).
  \item \celdaok{Modularización correcta}: separación ingesta / transformación /
        análisis / modelado / orquestación con lógica reutilizable en \texttt{src/}.
  \item \celdaok{Modelo dimensional adecuado}: esquema estrella en DuckDB con
        tabla puente N:N para la relación investigador--institución.
  \item \celdaok{Evidencias trazables}: 79 CSVs y 56 PNG/HTML con convención de
        nombres consistente por sprint.
  \item \celdaok{Enfoque crítico alineado con la evaluación responsable de la
        ciencia} (patrón \emph{Pregunta $\rightarrow$ Hallazgo $\rightarrow$ Caveat}).
\end{itemize}

\subsection{Principales brechas detectadas}

\begin{enumerate}
  \item \celdanok{CI/CD rota}: \texttt{.github/workflows/etl\_update.yml} ejecuta
        \texttt{src.ingesta.main} y \texttt{src.transformacion.main}, módulos que no
        existen (los reales son \texttt{src.ingesta.minciencias}/\texttt{produccion} y
        \texttt{src/Transformacion.py}).
  \item \celdanok{Dockerfile desalineado}: ejecuta \texttt{app/streamlit\_app.py}
        (placeholder de 17 líneas) en lugar de \texttt{streamlit\_app.py} (dashboard real).
  \item \celdanok{Mojibake en el CSV canónico}: \texttt{investigadores\_consolidado.csv}
        presenta caracteres corruptos (p.~ej. \texttt{Física} por \emph{Física}),
        indicativo de un desajuste UTF-8/Latin-1 en la exportación.
  \item \celdanok{Testing insuficiente}: 1 test funcional; \texttt{pandera} declarado
        pero sin uso; sin contratos de esquema.
  \item \celdanok{Legacy no reproducible}: scripts R/Python con rutas hardcodeadas
        (\texttt{Downloads/...}) y entradas no versionadas.
  \item \celdanok{Versionado de datos incompleto}: \texttt{datos/raw} y
        \texttt{datos/processed} gitignored sin DVC; la ingesta de producción
        (\texttt{33dq-ab5a}, ~1.2\,GB) no es reproducible desde el repositorio.
  \item \celdanok{Componente ML vacío}: \texttt{models/} sin contenido; el clustering
        (k-prototypes) vive en R legacy fuera del pipeline.
  \item \celdanok{Duplicidad de artefactos}: variantes \texttt{matriz\_*},
        \texttt{matriz\_*\_ext} y \texttt{matriz\_*\_obs} sin documentación de diferencias.
  \item \celdanok{Inconsistencia documental}: atípicos de edad reportados como 10 en
        un docstring y 22 en el README (el conteo real es 22).
  \item \celdanok{Nomenclatura desactualizada}: referencias a Grupo7, ramas y repos
        antiguos en el README.
\end{enumerate}

\subsection{Recomendaciones de máximo impacto}

\begin{enumerate}
  \item Corregir el workflow de GitHub Actions y el Dockerfile (reproducibilidad de
        la cadena de despliegue).
  \item Normalizar la codificación del CSV consolidado y añadir un contrato de
        esquema (pandera) en ingesta y transformación.
  \item Ampliar la cobertura de tests a las funciones analíticas y de transformación.
  \item Migrar los análisis legacy (R) al pipeline Python en \texttt{src/} o
        documentarlos formalmente como material de referencia histórica.
  \item Complementar la inferencia estadística: intervalos de confianza en índices
        de concentración, modelos de supervivencia para retención, modelos nulos
        para redes y validación formal de clusters.
\end{enumerate}

\subsection{Nota sobre la estructura de la entrega}

Este documento es el \textbf{compendio integral} de la auditoría: contiene la
documentación completa por archivo (sección 5) y la línea de tiempo completa
(sección 8). Por diseño, ese mismo contenido aparece en los documentos
\emph{autónomos} de los criterios 3 y 6 (secciones 3 de cada uno), que pueden
entregarse por separado; la duplicación entre el compendio y el documento de
criterio es intencional. Los criterios entre sí \textbf{no} repiten contenido:
cada documento de criterio desarrolla solo el suyo y remite aquí para el resto.
